% !TEX program = xelatex
\documentclass[11pt,a4paper]{article}
\usepackage[a4paper,margin=2.25cm]{geometry}
\usepackage{fontspec}
\usepackage{polyglossia}
\setmainlanguage{english}
\setmainfont{Times New Roman}
\usepackage{microtype}
\usepackage{amsmath,amssymb,booktabs}
\usepackage{hyperref}
\hypersetup{colorlinks=true,linkcolor=black,urlcolor=black}
\setlength{\parindent}{0pt}
\setlength{\parskip}{0.55em}
\renewcommand{\arraystretch}{1.12}
\title{\vspace{-1.2em}\textbf{Likelihood-weighted Candidate Semantic Agreement: a Paired RAGTruth QA Baseline}}
\author{}
\date{}
\begin{document}
\maketitle
\vspace{-1.7em}
\begin{abstract}
We evaluate a non-graph, response-conditioned hallucination baseline on the exact intersection of a completed semantic-sampling run and archive-verified HalluGraph scores.  The construct measures semantic agreement between the labelled historical answer and likely sampled answers to the same context and question.  This is a cache-only paired retrospective comparison, not a new Gemini-generation experiment or a factual-verification method.
\end{abstract}

\section{Construct}
For source prompt $x$, historical candidate answer $A$, and $15$ cached Gemini samples $s_i$, let $p_i$ denote normalized likelihood mass derived from selected-token mean log likelihood.  Candidate agreement and disagreement are
\begin{equation}
 M(A\mid x)=\sum_{i=1}^{15}p_i\,\mathbf{1}[A\equiv s_i],\qquad D(A\mid x)=1-M(A\mid x).
\end{equation}
Here $A\equiv s_i$ is the same bidirectional non-contradictory NLI relation used by the semantic-entropy implementation: neither direction may be a contradiction and the two directions may not both be neutral.  Larger $D$ is interpreted as less semantic support for the particular labelled candidate under likely sampled answers.  Unlike prompt-only semantic entropy, this construct explicitly scores $A$.

\section{Paired cache-only protocol}
The analysis uses the exact response-ID intersection between the completed 559-source semantic-entropy run and the R12 graph-score archive: 496 responses, comprising 405 paired training responses and 91 held-out responses (55 hallucinated and 36 factual).  The 63 completed entropy responses without a graph score are excluded rather than compared against a different denominator.

All 7,440 required semantic samples were loaded from the pinned compatible cache.  Consequently this evaluation made zero Gemini calls.  It performed 14880 local directional NLI evaluations to fill a separate content-addressed candidate-comparison cache.  Thresholds for every method were selected only on the 405 paired training responses.  A subsequent cache-only replay made zero Gemini calls and zero NLI evaluations, reproduced score records byte-for-byte, and left the candidate-comparison cache inventory unchanged.

\section{Held-out paired results}
\begin{table}[h]
\centering
\caption{Common $n=91$ held-out denominator.  Thresholds are fixed on the paired training split only.  Intervals are nonparametric bootstrap 95\% intervals.}
\small
\begin{tabular}{@{}lcccc@{}}
\toprule
Method & $\theta$ & ROC-AUC (95\% CI) & Precision / Recall / F1 (F1 95\% CI) \\
\midrule
candidate-agreement & -0.000 & 0.588 [0.458, 0.699] & 0.604 / 1.000 / 0.753 [0.662, 0.826] \\
strict & 0.168 & 0.763 [0.659, 0.852] & 0.640 / 1.000 / 0.780 [0.697, 0.849] \\
support & 0.203 & 0.751 [0.647, 0.853] & 0.719 / 0.745 / 0.732 [0.634, 0.824] \\
support-critical & 0.187 & 0.876 [0.804, 0.938] & 0.732 / 0.945 / 0.825 [0.746, 0.889] \\
\bottomrule
\end{tabular}
\end{table}

Relative to support-critical on the same held-out IDs, candidate-agreement has $\Delta$ROC-AUC = -0.288 (paired bootstrap 95\% interval [-0.422, -0.141]), and $\Delta$F1 = -0.072 ([-0.127, -0.017]).  These are paired uncertainty summaries on one fixed manifest; they are not an independent replication or a formal significance claim.

\section{Interpretation and limitations}
Candidate agreement is a meaningful non-graph response-level baseline because it tests the labelled answer rather than generator uncertainty alone.  In this setting, however, it is much weaker at ranking hallucinated answers than support-critical (ROC-AUC 0.588 versus 0.876).  Its train-selected threshold is effectively below the observed score range, yielding an all-positive held-out decision rule: recall is 1.000 but precision is limited by the 55/91 positive prevalence.  Thus its apparently moderate F1 should not be mistaken for discriminative evidence.

The construct is not direct factual verification.  A generator may consistently repeat a false claim, producing agreement with a hallucinated candidate; conversely, a correct but uncommon answer may obtain little likelihood mass.  It also inherits the sampled generator distribution and local NLI model.  Strict, support, and support-critical remain fundamentally different graph/evidence-based mechanisms.  This comparison establishes a family-diverse baseline on a shared denominator, not a replacement for evidence-grounded verification.

\end{document}
